\section{Mô hình bài toán}

\subsection{Biểu diễn bàn chơi và trạng thái}
Mỗi ô được định danh bằng tọa độ nguyên $(r,c)$. Bàn được nạp từ file JSON của màn chơi, gồm lưới ký tự (\texttt{\#} sàn, \texttt{.} ô trống, \texttt{S} xuất phát, \texttt{G} đích, \texttt{F} ô dễ vỡ, \texttt{B} cầu, \texttt{o}/\texttt{x}/\texttt{P} công tắc mềm/nặng/tách), danh sách \texttt{bridges} và danh sách \texttt{switches} mô tả liên kết giữa công tắc và cầu. Hàm \texttt{load\_board()} kiểm tra tính hợp lệ của màn trước khi Engine, GUI hay solver sử dụng. Ô cầu chỉ được xem là có thể đi khi trạng thái cầu tương ứng đang mở; ngoài bàn, ô trống và cầu đóng đều không đỡ khối, còn ô dễ vỡ không đỡ khối đang đứng.

Ở chế độ khối nguyên, trạng thái được biểu diễn bởi
\[
 s=(r,c,o,\mathbf{b}),
\]
trong đó $(r,c)$ là góc trên-trái của vùng chiếm chỗ, $o\in\{S,H,V\}$ lần lượt là đứng, nằm ngang và nằm dọc; $\mathbf{b}\in\{0,1\}^{B}$ là vector bit biểu diễn trạng thái đóng/mở của từng ô cầu, theo thứ tự khai báo trong \texttt{bridges}. Các ô khối chiếm là
\[
 C(S,r,c)=\{(r,c)\},\quad C(H,r,c)=\{(r,c),(r,c+1)\},\quad
 C(V,r,c)=\{(r,c),(r+1,c)\}.
\]
Trạng thái không cần lưu riêng trạng thái công tắc. Theo bảng chuyển ở mục sau, tập ô khối chiếm sau mỗi lượt lăn không giao với tập ô trước đó, nên mỗi lần khối đè lên công tắc luôn là một lần nhấn mới; khi khối đã tách, chỉ khối đơn vừa di chuyển mới kích hoạt công tắc. Vì vậy mọi thông tin ảnh hưởng tới các bước tiếp theo đã nằm trong $\mathbf{b}$. Hai trạng thái có cùng vị trí nhưng khác hướng hoặc khác vector cầu là hai nút khác nhau trong đồ thị tìm kiếm.

Khi công tắc tách được kích hoạt, trạng thái mở rộng thành
\[
 s=(p_1,p_2,k,\mathbf{b}),
\]
trong đó $p_1,p_2$ là vị trí hai khối đơn $1\times1\times1$ và $k\in\{1,2\}$ là khối đang được điều khiển. Trong mã nguồn, lớp \texttt{State} biểu diễn trường hợp này bằng \texttt{split\_mode = True}, \texttt{cube1\_position}, \texttt{cube2\_position}, \texttt{active\_cube}; \texttt{block\_position} luôn bằng vị trí khối đang điều khiển và \texttt{orientation} luôn là $S$. Ngay khi hai khối đơn nằm ở hai ô kề nhau, Engine tự động ghép chúng lại thành khối $1\times1\times2$ nằm ngang hoặc nằm dọc, với neo là ô trên-trái. Hai trạng thái tách có cùng vị trí nhưng khác khối đang điều khiển cũng là hai nút khác nhau.

\subsection{Hàm chuyển trạng thái}
Mỗi hành động thuộc tập $\{U,D,L,R,\mathit{SW}\}$, trong đó $\mathit{SW}$ (\texttt{SWITCH\_CUBE}) đổi khối đơn đang điều khiển và chỉ hợp lệ khi khối đang tách. Bảng dưới dùng tọa độ vùng chiếm chỗ đã định nghĩa ở trên; trạng thái mới chỉ được chấp nhận sau khi kiểm tra mọi ô mà khối chiếm.

\begin{table}[h!]
    \centering
    \caption{Chuyển hướng và tọa độ neo của khối nguyên}
    \label{tab:orientation_transitions}
    \begin{tabularx}{\textwidth}{c X X X X}
        \toprule
        \textbf{Hướng hiện tại} & \textbf{Lên ($U$)} & \textbf{Xuống ($D$)} & \textbf{Trái ($L$)} & \textbf{Phải ($R$)} \\
        \midrule
        Đứng $S(r,c)$ & $V(r-2,c)$ & $V(r+1,c)$ & $H(r,c-2)$ & $H(r,c+1)$ \\
        Ngang $H(r,c)$ & $H(r-1,c)$ & $H(r+1,c)$ & $S(r,c-1)$ & $S(r,c+2)$ \\
        Dọc $V(r,c)$ & $S(r-1,c)$ & $S(r+2,c)$ & $V(r,c-1)$ & $V(r,c+1)$ \\
        \bottomrule
    \end{tabularx}
\end{table}

Khi khối đang tách, hành động hướng chỉ dịch khối đơn đang điều khiển đúng một ô; nếu ô đích trùng vị trí khối đơn còn lại thì hành động không hợp lệ.

\texttt{Engine.transition()} xử lý một hành động theo thứ tự: (1) áp dụng bảng chuyển; (2) kiểm tra mọi ô mà khối chiếm theo vector cầu hiện tại: ô ngoài bàn, ô trống, cầu đóng làm khối rơi, khối nguyên đang đứng trên ô dễ vỡ cũng rơi, còn khối nằm được phép đi qua ô này; (3) áp dụng hiệu ứng ô: công tắc mềm kích hoạt khi bất kỳ phần nào của khối đè lên, công tắc nặng và công tắc tách chỉ kích hoạt khi khối nguyên đứng trên đó, khối đơn chỉ kích hoạt được công tắc mềm; mỗi công tắc cập nhật các cầu liên kết theo chế độ \texttt{toggle}, \texttt{permanent\_open} hoặc \texttt{permanent\_closed}; công tắc tách đưa hai khối đơn tới hai ô \texttt{split\_cells} đã khai báo và giao quyền điều khiển cho khối 1; sau cùng thử ghép khối; (4) nếu bước (3) làm thay đổi trạng thái, kiểm tra nền lần nữa, vì công tắc có thể vừa đóng cầu ngay dưới khối hoặc đưa khối đơn tới ô không có nền. Kết quả trả về là \texttt{OK} kèm trạng thái mới, \texttt{FALL} kèm lý do (\texttt{off\_board}, \texttt{void}, \texttt{bridge\_closed}, \texttt{fragile\_standing}), hoặc \texttt{ILLEGAL}. Điều kiện đích (\texttt{Engine.is\_goal()}) là khối không tách, đang đứng và nằm đúng ô \texttt{G}; khối nằm đè lên ô đích hay khối đơn trên ô đích đều không thắng.

\section{Thiết kế thuật toán}

\subsection{Bốn phương pháp tìm kiếm}
\begin{itemize}
    \item \textbf{BFS:} Dùng hàng đợi FIFO, mở rộng theo từng lớp số lượt đi. Với mỗi trạng thái, thêm các trạng thái kế tiếp hợp lệ chưa thăm. BFS tìm được lời giải có ít lượt nhất khi mọi lượt có cùng chi phí; đổi lại, số trạng thái lưu trong hàng đợi có thể tăng nhanh.
    \item \textbf{DFS:} Dùng ngăn xếp LIFO để đi sâu theo một nhánh trước khi quay lui. Nhóm dùng DFS dạng tìm kiếm trên đồ thị với tập trạng thái đã thăm để ngăn vòng lặp vô hạn, không dùng IDS và không giới hạn độ sâu. DFS có thể tìm được lời giải nhưng không đảm bảo lời giải ngắn nhất.
    \item \textbf{UCS:} Dùng hàng đợi ưu tiên theo chi phí đường đi $g(s)$ và luôn mở rộng trạng thái có chi phí nhỏ nhất. Với cạnh có chi phí không âm, UCS trả về lời giải có tổng chi phí nhỏ nhất. Chi phí đề xuất dưới đây không đồng nhất, do đó UCS không bị rút gọn thành BFS.
    \item \textbf{$A^*$:} Ưu tiên trạng thái có $f(s)=g(s)+h(s)$, trong đó $g$ là chi phí đã đi và $h$ là cận dưới chi phí còn lại. Với heuristic admissible và chi phí không âm, $A^*$ tìm được lời giải tối ưu; khi heuristic hữu ích, thuật toán thường mở rộng ít nút hơn UCS, nhưng vẫn có thể tốn nhiều bộ nhớ \cite{hart1968formal,russell2020aima}.
\end{itemize}

Mỗi thuật toán lưu nút cha và hành động đã dùng để dựng lại đường đi sau khi tìm thấy đích. BFS đánh dấu một trạng thái là đã phát hiện khi đưa vào frontier, còn DFS đánh dấu khi lấy trạng thái ra khỏi stack. Với UCS/$A^*$, lưu chi phí tốt nhất đã biết \texttt{best\_cost} và chỉ cập nhật khi tìm thấy đường đi rẻ hơn; bản ghi cũ trong hàng đợi ưu tiên bị bỏ qua khi lấy ra (xóa trễ), và một bộ đếm tăng dần được dùng để phá hòa theo thứ tự chèn. Vì \texttt{State} là dataclass bất biến và băm được, khóa băm bao gồm toàn bộ trạng thái, kể cả vector cầu và cấu hình tách khối.

\subsection{Cơ chế dùng chung của các solver tìm kiếm}

\subsubsection{Lớp nền và trạng thái đầu}
Các solver BFS và DFS dùng chung lớp nền \texttt{BaseSolver}. Lớp này kiểm tra trạng thái đầu, gọi phương thức tìm kiếm, truy vết lời giải và đóng gói \texttt{SearchResult}. Mỗi lớp con chỉ hiện thực \texttt{\_search()} với chiến lược và frontier riêng. UCS và $A^*$ được viết thành hàm \texttt{solve()} độc lập trong \texttt{solvers/ucs.py} và \texttt{solvers/astar.py}, nhưng dùng lại các hàm \texttt{reconstruct\_path()} và \texttt{finalize()} của \texttt{solvers/base\_solver.py}.

Trước khi bắt đầu tìm kiếm, \texttt{BaseSolver.solve()} kiểm tra trạng thái đầu. Nếu trạng thái đầu đã là đích, đường đi rỗng, cha là \texttt{None} và số trạng thái mở rộng bằng không. $A^*$ tự kiểm tra trường hợp này trước khi dựng heuristic; UCS lấy trạng thái đầu ra khỏi hàng đợi và kiểm tra đích trước khi mở rộng, nên cũng cho cùng kết quả.

Cả bốn solver sinh trạng thái kế tiếp qua \texttt{Engine.successors()}. Hàm này duyệt hành động theo thứ tự cố định \texttt{UP}, \texttt{DOWN}, \texttt{LEFT}, \texttt{RIGHT}, \texttt{SWITCH\_CUBE}, gọi \texttt{transition()} đúng một lần cho mỗi hành động và chỉ trả về các cặp (hành động, trạng thái mới) có kết quả \texttt{OK}.

\subsubsection{Truy vết và xác nhận đường đi}
Các solver lưu thông tin cần thiết để truy vết, thay vì giữ toàn bộ đường đi cho mọi nút. Khi gặp đích, \texttt{reconstruct\_path()} lần ngược từ đích về gốc, thu thập hành động rồi đảo danh sách để nhận đúng thứ tự thực hiện. Quá trình truy vết cũng kiểm tra yêu cầu hủy.

Sau khi khôi phục, \texttt{path\_cost()} phát lại từng hành động qua Engine và tính tổng chi phí bằng \texttt{Engine.step\_cost()}:
\begin{lstlisting}[language=Python]
result = engine.transition(state, action)
if result.status is not TransitionStatus.OK or result.state is None:
    raise RuntimeError("Path contains an invalid action")

total += engine.step_cost(state, action, result.state)
state = result.state
\end{lstlisting}
Bước xác nhận bảo đảm đường đi trả về còn hợp lệ theo Engine và cung cấp tổng chi phí. Việc tính được chi phí không tự đảm bảo lời giải tối ưu; tính tối ưu phụ thuộc thuật toán và các điều kiện của heuristic nếu có.

\subsubsection{Kết quả và cơ chế hủy}
\texttt{SearchResult} lưu trạng thái giải được, đường đi, tổng chi phí, thời gian, bộ nhớ, số trạng thái mở rộng, độ dài lời giải và thời gian tiền xử lý. Thời gian đo bằng \texttt{time.perf\_counter()}; bộ nhớ là đỉnh cấp phát đo bằng \texttt{tracemalloc}, tính bằng MB. Với $A^*$, thời gian dựng heuristic được tách sang \texttt{preprocessing\_time} và không cộng vào thời gian tìm kiếm, còn bộ nhớ đỉnh vẫn bao gồm bảng heuristic. Các solver kiểm tra yêu cầu hủy trong vòng lặp; cơ chế này chỉ phục vụ giới hạn thời gian của GUI và không làm thay đổi nguyên tắc duyệt cốt lõi.

\subsection{Triển khai Breadth-First Search (BFS)}

\subsubsection{Ý tưởng và cấu trúc dữ liệu}
Breadth-First Search xử lý trạng thái theo từng mức độ sâu tính từ trạng thái ban đầu. Vì BFS xem mọi cạnh có cùng trọng số, thuật toán tìm được đường đi có số hành động nhỏ nhất. Triển khai nằm trong lớp \texttt{BFSSolver} tại \texttt{solvers/bfs.py}.

BFS dùng hàng đợi FIFO \texttt{collections.deque}, bản đồ \texttt{parent} để đánh dấu trạng thái đã phát hiện và lưu thông tin truy vết, cùng biến đếm trạng thái mở rộng:
\begin{lstlisting}[language=Python]
parent: ParentMap = {initial_state: None}
queue: deque[State] = deque([initial_state])
expanded_nodes = 0
\end{lstlisting}
Kiểu dữ liệu \texttt{ParentMap} trong \texttt{solvers/base\_solver.py}:
\begin{lstlisting}[language=Python]
ParentMap = dict[State, tuple[State, Action] | None]
\end{lstlisting}
Giá trị \texttt{None} biểu thị trạng thái gốc; giá trị còn lại lưu trạng thái cha và hành động đã sinh ra trạng thái hiện tại.

\subsubsection{Luồng xử lý}
Phương thức tìm kiếm chính là \texttt{BFSSolver.\_search()}. Mỗi vòng lặp kiểm tra yêu cầu hủy rồi lấy trạng thái ở đầu hàng đợi. Nếu hàng đợi rỗng, thuật toán kết thúc mà không tìm thấy đích.
\begin{lstlisting}[language=Python]
while queue:
    check_cancelled()
    current = queue.popleft()
\end{lstlisting}
BFS kiểm tra điều kiện đích ngay sau khi lấy trạng thái ra khỏi hàng đợi, không phải ngay khi phát hiện successor:
\begin{lstlisting}[language=Python]
if self.engine.is_goal(current):
    return current, parent, expanded_nodes
\end{lstlisting}
Nếu chưa phải đích, BFS tăng số trạng thái mở rộng, duyệt các successor hợp lệ do Engine cung cấp và enqueue trạng thái mới chưa được phát hiện:
\begin{lstlisting}[language=Python]
expanded_nodes += 1
for action, nxt in self.engine.successors(current):
    check_cancelled()
    if nxt not in parent:
        parent[nxt] = (current, action)
        queue.append(nxt)
\end{lstlisting}
Khi frontier cạn, hàm trả về \texttt{None}, bản đồ cha và số trạng thái đã mở rộng.

\subsubsection{Mã giả BFS theo triển khai}
\begin{lstlisting}
BFS(initial_state):
    parent[initial_state] = None
    queue = deque([initial_state])
    expanded_nodes = 0
    while queue is not empty:
        check_cancelled()
        current = queue.popleft()
        if is_goal(current):
            return current, parent, expanded_nodes
        expanded_nodes += 1
        for (action, next_state) in successors(current):
            check_cancelled()
            if next_state not in parent:
                parent[next_state] = (current, action)
                queue.append(next_state)
    return None, parent, expanded_nodes
\end{lstlisting}
Mã giả tóm lược logic trong solver, không thay thế mã nguồn chính thức.

\subsubsection{Đặc điểm đáng chú ý của BFS}
BFS ghi trạng thái vào \texttt{parent} ngay khi enqueue. Vì vậy, mỗi trạng thái chỉ được đưa vào hàng đợi một lần và đường đi đầu tiên phát hiện trạng thái được giữ để truy vết. \texttt{Engine.successors()} trả về hành động theo thứ tự cố định nên thứ tự duyệt giữa các trạng thái cùng độ sâu ổn định với cùng đầu vào. Khi có nhiều lời giải ngắn nhất, thứ tự successor quyết định lời giải nào được chọn nhưng không đổi tính tối ưu theo số hành động.

BFS không đọc \texttt{step\_cost()} trong quá trình tìm kiếm; tổng chi phí chỉ được tính khi phát lại đường đi. Do đó lời giải ít bước nhất của BFS có thể đắt hơn lời giải của UCS (xem mục \ref{subsec:bfs_dfs_props}).

\subsection{Triển khai Depth-First Search (DFS)}

\subsubsection{Ý tưởng và cấu trúc dữ liệu DFS}
Depth-First Search đi sâu theo một nhánh trước và chuyển sang nhánh khác khi không thể tiếp tục. Trong dự án, DFS dùng ngăn xếp tường minh và được triển khai trong lớp \texttt{DFSSolver} tại \texttt{solvers/dfs.py}. Nhóm chọn DFS dạng tìm kiếm trên đồ thị, không dùng Iterative Deepening Search.

Mỗi phần tử stack gồm trạng thái, trạng thái cha ứng viên và hành động dẫn tới trạng thái đó:
\begin{lstlisting}[language=Python]
stack: list[tuple[State, State | None, Action | None]] = [
    (initial_state, None, None)
]
visited: set[State] = set()
parent: ParentMap = {}
expanded_nodes = 0
\end{lstlisting}
Danh sách Python với \texttt{append()} và \texttt{pop()} ở cuối thực hiện nguyên tắc LIFO. DFS dùng \texttt{visited} để lưu trạng thái đã xử lý, còn \texttt{parent} lưu thông tin truy vết đã được chấp nhận.

\subsubsection{Luồng xử lý DFS}
Mỗi vòng lặp kiểm tra yêu cầu hủy rồi lấy phần tử trên cùng của stack. Vì trạng thái được đánh dấu khi pop, một trạng thái chưa thăm có thể được thêm vào stack nhiều lần; bản sao lấy ra sau sẽ bị bỏ qua.
\begin{lstlisting}[language=Python]
while stack:
    check_cancelled()
    current, previous, action_from_previous = stack.pop()
    if current in visited:
        continue
\end{lstlisting}
Trạng thái mới được đánh dấu và ghi nhận cha; trạng thái gốc có cha là \texttt{None}:
\begin{lstlisting}[language=Python]
visited.add(current)
parent[current] = (
    None
    if previous is None or action_from_previous is None
    else (previous, action_from_previous)
)
\end{lstlisting}
Sau khi ghi nhận trạng thái, DFS kiểm tra đích. Nếu chưa phải đích, thuật toán tăng bộ đếm rồi lưu successor hợp lệ chưa được thăm:
\begin{lstlisting}[language=Python]
if self.engine.is_goal(current):
    return current, parent, expanded_nodes

expanded_nodes += 1
successors: list[tuple[State, Action]] = []
for action, nxt in self.engine.successors(current):
    check_cancelled()
    if nxt not in visited:
        successors.append((nxt, action))
\end{lstlisting}
Danh sách successor được đảo trước khi đưa vào stack:
\begin{lstlisting}[language=Python]
for next_state, action in reversed(successors):
    stack.append((next_state, current, action))
\end{lstlisting}
Engine trả hành động theo thứ tự \texttt{UP}, \texttt{DOWN}, \texttt{LEFT}, \texttt{RIGHT}, \texttt{SWITCH\_CUBE}. Đảo thứ tự trước khi push khiến successor đầu danh sách được pop trước. Khi stack rỗng mà chưa gặp đích, DFS trả về \texttt{None}, \texttt{parent} và \texttt{expanded\_nodes}.

\subsubsection{Mã giả DFS theo triển khai}
\begin{lstlisting}
DFS(initial_state):
    stack = [(initial_state, None, None)]
    visited = set()
    parent = {}
    expanded_nodes = 0
    while stack is not empty:
        check_cancelled()
        current, previous, action = stack.pop()
        if current in visited:
            continue
        visited.add(current)
        parent[current] = None if previous is None else (previous, action)
        if is_goal(current):
            return current, parent, expanded_nodes
        expanded_nodes += 1
        children = []
        for (action, next_state) in successors(current):
            check_cancelled()
            if next_state not in visited:
                children.append((next_state, action))
        for next_state, action in reversed(children):
            stack.append((next_state, current, action))
    return None, parent, expanded_nodes
\end{lstlisting}
Mã giả lược bỏ cú pháp phụ để làm rõ cấu trúc của \texttt{DFSSolver.\_search()}.

\subsubsection{Đặc điểm đáng chú ý của DFS}
Đánh dấu khi pop có thể làm cùng trạng thái xuất hiện nhiều lần trong stack. Khi một bản sao được xử lý, các bản sao còn lại bị bỏ qua. Thông tin cha được ghi lần đầu trạng thái được chấp nhận xử lý; các phần tử trùng lặp không ghi đè đường truy vết.

DFS không ưu tiên độ sâu nhỏ nhất. Vì vậy, lời giải đầu tiên tìm được không nhất thiết ngắn nhất hay rẻ nhất. Thứ tự successor quyết định nhánh được đi sâu trước nhưng không tạo bảo đảm tối ưu.

\subsection{Tính chất và độ phức tạp của BFS/DFS}
\label{subsec:bfs_dfs_props}

\subsubsection{Tính đầy đủ và tối ưu}
Tính đầy đủ là khả năng tìm được lời giải nếu lời giải tồn tại trong không gian trạng thái đang xét. Với không gian hữu hạn, mô hình trạng thái đầy đủ và thuật toán được chạy đến khi hoàn tất, BFS duyệt các lớp có thể tiếp cận; DFS dùng \texttt{visited} để tránh mở rộng lặp lại trạng thái đã xử lý. Nếu quá trình bị hủy hoặc dừng sớm, không thể khẳng định đã tìm hết không gian.

BFS tìm đường có số cạnh nhỏ nhất trong đồ thị có trọng số cạnh bằng nhau; ở đây mỗi cạnh là một hành động hợp lệ. Điều đó không đảm bảo tổng \texttt{step\_cost} nhỏ nhất vì chi phí thay đổi theo tư thế khối và loại hành động: trên bộ thực nghiệm, có 7 case mà lời giải BFS có cùng số bước với UCS nhưng tổng chi phí lớn hơn. DFS trả về lời giải đầu tiên gặp theo thứ tự duyệt, không bảo đảm tối thiểu số hành động hay tổng chi phí.

\subsubsection{Độ phức tạp và giới hạn}
Gọi $V$ là số trạng thái có thể tiếp cận và $E$ là số chuyển trạng thái hợp lệ. Với $N$ ô không phải ô trống và $B$ ô cầu, $V$ bị chặn trên bởi khoảng $(3N + 2N(N-1))\cdot 2^{B}$: ba hướng cho khối nguyên, cặp vị trí có thứ tự cùng khối đang điều khiển cho khối đã tách, nhân với mọi cấu hình cầu. Trên thực tế chỉ một phần nhỏ trong số này đạt được. Mỗi trạng thái có tối đa 5 hành động nên $E\le 5V$. Với tìm kiếm đồ thị chuẩn, BFS có thời gian $O(V+E)$ và bộ nhớ $O(V)$. DFS với tập \texttt{visited} và lưu cha cũng thường được mô tả bằng $O(V+E)$ thời gian và $O(V)$ bộ nhớ.

Trong biến thể DFS của dự án, trạng thái được đánh dấu khi pop khỏi stack. Vì vậy, trạng thái chưa thăm có thể được đưa vào stack từ nhiều nhánh; số phần tử stack trung gian có thể lớn hơn số trạng thái duy nhất, nhưng không vượt quá $E$ vì mỗi trạng thái chỉ được mở rộng một lần.

Theo hệ số phân nhánh $b$ và độ sâu lời giải $d$, BFS thường được mô tả với chi phí tăng theo cấp số mũ của $d$, xấp xỉ $O(b^d)$; DFS thường được mô tả là $O(b^m)$ với $m$ là độ sâu tối đa. Đây là mô hình phân tích cây tìm kiếm, khác với biểu diễn theo $V,E$ và cần nêu giả định khi sử dụng.

\subsection{Chi phí UCS và heuristic $A^*$}

\subsubsection{Hàm chi phí}
Để UCS tối ưu một tiêu chí khác số bước đơn thuần, \texttt{Engine.step\_cost()} đặt chi phí một lượt theo tư thế của trạng thái nguồn:
\begin{equation}
    c(s,a,s')=
    \begin{cases}
        0.1 & \text{nếu } a=\mathit{SW},\\
        1.5 & \text{nếu } a\neq\mathit{SW} \text{ và } o(s)=S,\\
        1.0 & \text{nếu } a\neq\mathit{SW} \text{ và } o(s)\in\{H,V\}.
    \end{cases}
    \label{eq:ucs_cost}
\end{equation}
Các giá trị được khai báo trong \texttt{CostConfig} (\texttt{stand\_roll}, \texttt{lie\_roll}, \texttt{switch\_cube}). Lăn từ tư thế đứng làm khối lật qua cạnh và dịch tọa độ neo hai ô, đồng thời đưa khối về tư thế nằm; lăn khi đang nằm chỉ dịch một ô. Chi phí cao hơn cho lượt lăn từ tư thế đứng khuyến khích lời giải giữ khối ở tư thế nằm ổn định và hạn chế các lần dựng khối không cần thiết. Đổi khối điều khiển không làm khối di chuyển nên có chi phí rất nhỏ nhưng vẫn dương, tránh chuỗi đổi qua lại vô ích. Vì khối đơn luôn có hướng $S$, mỗi bước di chuyển khối đơn có chi phí 1.5. Mọi chi phí đều hữu hạn và dương, nên UCS và $A^*$ hoạt động đúng; vì các chi phí khác nhau, UCS không trùng với BFS.

\subsubsection{Heuristic $h_2$: khoảng cách trên đồ thị nới lỏng}
$A^*$ mặc định dùng heuristic $h_2$ (\texttt{solvers/heuristics/h2.py}), được xây bằng cách trừu tượng hóa bài toán rồi tính chính xác khoảng cách trên bài toán đã trừu tượng.

\textbf{Ánh xạ trừu tượng.} Với khối nguyên, $\alpha(s)=(r,c,o)$; với khối đã tách, $\alpha(s)=\{p_1,p_2\}$ là cặp vị trí không thứ tự. Ánh xạ bỏ vector cầu $\mathbf{b}$ và khối đang điều khiển $k$.

\textbf{Nền đỡ nới lỏng.} Gọi $R$ là tập mọi ô không phải ô trống. Trong bài toán trừu tượng, ô dễ vỡ được xem như sàn thường, cầu luôn mở, và công tắc mềm/nặng không có hiệu ứng.

\textbf{Đồ thị trừu tượng $G_\alpha$.} Đỉnh gồm mọi tư thế $(r,c,o)$ có $C(o,r,c)\subseteq R$, cùng các cặp khối đơn đạt được từ \texttt{split\_cells} của công tắc tách. Cạnh gồm:
\begin{itemize}
    \item lượt lăn theo bảng \ref{tab:orientation_transitions} tới tư thế nằm trong $R$, trọng số 1.5 nếu tư thế nguồn đứng và 1.0 nếu đang nằm;
    \item lượt lăn kết thúc ở tư thế đứng trên công tắc tách được nối thẳng tới cặp \texttt{split\_cells} (hoặc tư thế đã ghép nếu hai ô kề nhau), trọng số 1.5;
    \item từ một cặp, di chuyển \emph{bất kỳ} khối đơn nào một ô vào $R$ mà không trùng khối còn lại, trọng số 1.5; nếu hai khối kề nhau thì ghép thành tư thế nằm, neo ở ô trên-trái.
\end{itemize}
Hành động $\mathit{SW}$ không có cạnh vì nó không đổi $\alpha(s)$. Heuristic được định nghĩa là
\begin{equation}
    h_2(s)=d_\alpha\bigl(\alpha(s)\bigr),
    \label{eq:astar_heuristic}
\end{equation}
với $d_\alpha(x)$ là khoảng cách ngắn nhất từ $x$ tới đỉnh $(r_g,c_g,S)$ trên $G_\alpha$, và $d_\alpha(x)=\infty$ nếu không tới được. Bảng $d_\alpha$ được tính một lần cho mỗi màn bằng Dijkstra ngược từ đỉnh đích; trong lúc tìm kiếm, mỗi lần đánh giá $h_2$ chỉ là một phép tra bảng $O(1)$. Thời gian dựng bảng được ghi vào \texttt{preprocessing\_time}.

\textbf{Admissible và consistent.} Xét một chuyển trạng thái hợp lệ $s\xrightarrow{a}s'$. Nếu $a=\mathit{SW}$ thì $\alpha(s)=\alpha(s')$, nên $h_2(s)=h_2(s')\le 0.1+h_2(s')$. Ngược lại, mọi ô mà khối chiếm sau lượt đi đều được nền đỡ trong bài toán thật, nên đều không phải ô trống và thuộc $R$; công tắc tách đưa khối tới đúng cặp \texttt{split\_cells} như cạnh tương ứng trong $G_\alpha$; khối đơn đang điều khiển là một trong hai khối mà $G_\alpha$ cho phép di chuyển; quy tắc ghép khối giống nhau ở hai phía. Do đó $G_\alpha$ luôn có cạnh $\alpha(s)\to\alpha(s')$ với trọng số đúng bằng $c(s,a,s')$. Các luật bị nới lỏng (ô dễ vỡ, cầu đóng, va chạm của khối đơn với ô không có nền) chỉ làm bài toán thật có \emph{ít} cạnh hơn. Theo bất đẳng thức tam giác của đường đi ngắn nhất,
\[
 h_2(s)\le c(s,a,s')+h_2(s'),
\]
tức $h_2$ consistent. Trạng thái đích có $\alpha=(r_g,c_g,S)$ nên $h_2=0$. Cộng dồn bất đẳng thức trên dọc một đường đi tối ưu tới đích cho $h_2(s)\le h^*(s)$, tức $h_2$ admissible với chính hàm chi phí \eqref{eq:ucs_cost}. Hệ quả: $h_2(s)=\infty$ kéo theo $h^*(s)=\infty$, nên $A^*$ không đưa các trạng thái này vào frontier mà không mất lời giải. Trọng số cạnh phải tính theo tư thế nguồn như trên; nếu chỉ đếm số bước rồi nhân với chi phí nhỏ nhất 1.0, heuristic vẫn an toàn nhưng yếu hơn.

\subsubsection{Heuristic cơ sở $h_1$}
Để đánh giá lợi ích của $h_2$, nhóm cài thêm heuristic cơ sở (\texttt{solvers/heuristics/h1.py}):
\begin{equation}
    h_1(s)=\left\lceil\frac{|r-r_g|+|c-c_g|}{2}\right\rceil,
    \label{eq:h1_baseline}
\end{equation}
trong đó $(r,c)$ là góc trên-trái của vùng chiếm khối và $(r_g,c_g)$ là ô đích; $h_1=0$ khi khối đã tách. Một lượt lăn thay đổi tọa độ neo nhiều nhất hai ô theo khoảng cách Manhattan, nên $h_1$ giảm không quá 1 sau mỗi lượt lăn có chi phí ít nhất 1.0. Tuy nhiên, công tắc tách có thể đưa khối tới gần đích hơn khoảng cách Manhattan cho phép, và khi đó $h_1$ giảm đột ngột về 0 sau một cạnh có chi phí 1.5. Vì vậy $h_1$ không đảm bảo consistent, và cũng không có chứng minh admissible trên màn có công tắc tách. $h_1$ chỉ được dùng để so sánh, không phải heuristic mặc định của $A^*$.

\subsection{Kiểm chứng tính đúng đắn}
Repo chưa có bộ unit test tự động; tính đúng đắn được kiểm chứng bằng các cơ chế sau. Thứ nhất, \texttt{load\_board()} kiểm tra mỗi màn có đúng một ô xuất phát và một ô đích, và mọi ô cầu, ô công tắc trên lưới khớp với khai báo trong \texttt{bridges}/\texttt{switches}. Thứ hai, mọi lời giải đều được phát lại qua \texttt{Engine.transition()} trong \texttt{path\_cost()}; một hành động không hợp lệ sẽ làm solver báo lỗi. Thứ ba, trên 120 case cố định, cả bốn thuật toán giải được 480/480 lượt, BFS, UCS và $A^*$ cho cùng độ dài lời giải, và tổng chi phí của $A^*$ trùng với UCS ở mọi case. Thứ tư, \texttt{experiments/verify\_heuristic.py} duyệt toàn bộ 115.043 trạng thái đạt được và 379.858 cạnh hợp lệ của 120 case, tính $h^*$ chính xác bằng Dijkstra ngược từ các trạng thái đích, rồi kiểm tra từng trạng thái và từng cạnh: $h_2$ không có vi phạm admissible hay consistent nào; $h_1$ không vi phạm admissible trên bộ test này nhưng có 21 cạnh vi phạm consistent, đều nằm ở các màn có công tắc tách (GUI-008, 009, 010, 015, 017).
